\newpage

\section{\huge \textcolor{blue}{章节名}}

\subsection{\LARGE 子标题}

\subsubsection{\large 子子标题}

\hspace*{1em}
%带序号的公式
{\Large
\begin{equation}
a(x) y' + b(x) y = c(x)
\end{equation}
}
%无编号的公式
{\Large \[ (\mu(x) y)' = \mu(x) q(x) \]}



\begin{enumerate}
    \item 
    \[  A=\z[a^2,a^2+1\y)    \]
    
    \item 
    \[ B=\z[\frac{a^2}{\sqrt[3]{b}},\frac{a^2+1}{\sqrt[3]{b}}\y) \]
    
    \item
    \[ C=\z\{x\in\R: \frac{x^2}{x^2+1}>1 \y\}    \]
\end{enumerate}

\begin{definition}[有理数]
我们称$p/q$是一个有理数，当且仅当
\begin{enumerate}
    \item $p,q$是两个整数。
    \item $q$不等于0。
    \item $p,q$的最大公因数是1，或$p,q$互素。
\end{enumerate}

我们称由全体有理数构成的集合为有理数集，记作$\Q$。
\end{definition}

\begin{align*}
    &=  \\
    &=   \\
    &= \\
    &=      
\end{align*}

\[
\begin{vmatrix}
    y_1 & y_2 \\
    y'_1 & y'_2
\end{vmatrix}_{x_0}
\neq 0
\]
行列式


\begin{center}
    \begin{tabular}{c|c|c|c|c}
        $n$ & $x_n$ & $y_n$ & $h_n$ & $h f_n$ \\ \hline
        0   & 0     & 1     & $-1$   & $-0.1$   \\ \hline
        1   & 0.1   & 0.9   & $-0.80$ & $-0.08$  \\ \hline
        2   & 0.2   & \textcolor{red}{0.82}   &        &          \\
    \end{tabular}
\end{center}

\[\left\{
\begin{aligned}
   u &= u  \\
   v &= t - u
\end{aligned}
\right.\]

\textcolor{red}{not}

$\begin{cases}
\end{cases}$

\begin{figure}[htbp]
    \centering
\includegraphics[width=0.5\textwidth]{C:/Users/37/Desktop/Notes on Common Calculus Equations/Notes and pictures/cwf10.4.7.png}
    \caption{平移-阶跃-拉普拉斯变换}
    \label{fig:cwf10.4.7}
\end{figure}